% Fragmento integrable. Preambulo anfitrion: amsmath, amssymb, longtable, hyperref.
\section{Isolation of the limit and scaling of the global selector}
\label{hmtfg:fragmento:aislamiento-del-limite-y-escalado-del-selector-global}

\subsection{Indefinite continuation of the first root and exclusion of the intermediate interval}
\label{hmtfg:escalado:15}

The complete classification proved above preserves the family and the first-new-root selector. Ordered transport preserves the zeros and their minimum; the return word preserves the nonadic chronological reading.

\subsubsection{Classification and existence at every depth}
\label{hmtfg:escalado:15.1}

Write \(W_n=\tau_1\cdots\tau_{2^n-1}\), with \(\tau_j=(-1)^{v_2(j)}\), and \(a_*=\min\Lambda_\tau\). The existence and compactness of \(\Lambda_\tau\) are proved in the prolongation development, \ref{hmtfg:prolongacion:6.1}–\ref{hmtfg:prolongacion:6.2}; its minimum exists by compactness and order.

\textbf{Classification.} Every strictly unimodal map with a maximum, on an invariant interval, whose critical point has exact period \(2^n\) and whose itinerary \(K\) satisfies \(K<\tau\), has itinerary \(W_n0\).

The induction proceeds through the effective return. For period at least eight, comparison with \(\tau\), together with the intervals that would trap orbits with incompatible signs, forces the prefix \((+,-,+,+,+,-)\). This yields
\[
J_1=[d_2,d_4],\quad F(J_1)=[d_3,1],\quad
d_4<d_3,\quad F^2(J_1)=J_1.
\]
The return normalized by \(h(x)=d_2x\) has maximum one, half the period, and itinerary \(K'_j=-K_{2j}<\tau\). Order is preserved because
\[
O_{2q-1}(K)=-O_{q-1}(K'),\qquad
K_{2q}-\tau_{2q}=-(K'_q-\tau_q).
\]
The initial cases are \(+0\) and \(+-+0\); induction reconstructs \(W_n0\). For even maps, also \(0<d_4<-d_2<1\), so the return extends to \([-1,1]\) and coincides with \eqref{hmtfg:escalado:eq:6}.

Before \(a_*\), all itineraries are smaller than \(\tau\): the strict-comparison sets are open by the separation proved in \ref{hmtfg:prolongacion:6.1}; the interval preceding \(a_*\) is connected and begins with \(+^\infty<\tau\).

\textbf{Existence of each first root.} Suppose \(\lambda_{n-1}<a_*\) has been constructed, and let \(p=2^{n-1}\). The zeros of \(c_p\) in \([\lambda_{n-1},a_*]\) form a nonempty finite set: the function is analytic and \(c_p(a_*)\ne0\). Let \(z<a_*\) be its last zero. In \((z,a_*]\), \(c_p\) has sign \(\tau_p\). For \(g_\lambda=f_\lambda^p\),
\[
g_\lambda'(0)=0,\qquad
c_{2p}(\lambda)=c_p(\lambda)+O(c_p(\lambda)^2).
\]
Near \(z\), on the right, \(c_{2p}\) has sign \(\tau_p\); at \(a_*\), it has the opposite sign. There is a zero of \(c_{2p}\) in \((z,a_*)\), and its period is exactly \(2p\), since there \(c_p\ne0\). Finiteness of the zeros provides the first new root in \((\lambda_{n-1},a_*)\). The last auxiliary zero proves existence; the effective choice remains the first root. The base case is \(\lambda_1=1/u_0(1)<a_*\).

\subsubsection{Limit of the original sequence}
\label{hmtfg:escalado:15.2}

\textbf{Theorem.} The original selector exists at every level and satisfies
\begin{equation}
\boxed{K_{\lambda_n}=W_n0,\qquad
\lambda_n\nearrow a_*=\min\Lambda_\tau.}
\label{hmtfg:escalado:eq:50}
\end{equation}
\textbf{Proof.} The preceding existence result gives an increasing sequence bounded by \(a_*\). The classification fixes all its itineraries. Let \(b\le a_*\) be its limit. For each \(p\), sufficiently late terms preserve the opposite signs \(\tau_p,\tau_{2p}\). Taylor’s formula and the uniform bound \(B_p=16^p(16^p-1)/15\) give
\[
|c_{2p}(\lambda_n)-c_p(\lambda_n)|
\le\tfrac12 B_p|c_p(\lambda_n)|^2,\qquad
|c_p(\lambda_n)|>2/B_p.
\]
Passing to the limit preserves \(\operatorname{sign}c_p(b)=\tau_p\), with absolute value at least \(2/B_p>0\). For every \(p\), this implies \(b\in\Lambda_\tau\); minimality gives \(b=a_*\).

The result establishes indefinite existence, preservation of the symbolic cascade, and selection of its first limiting parameter for the original global rule. The HMT completion transports this sequence and its limit while preserving order and reading.

\subsubsection{Certificate excluding the intermediate interval}
\label{hmtfg:escalado:15.3}

The bridge program certifies
\begin{equation}
\Lambda_\tau\cap[\lambda_{8,\rm inf},1.874038]=\varnothing,
\label{hmtfg:escalado:eq:51}
\end{equation}
where
\[
\lambda_{8,\rm inf}
=1.874027744154450672897007573595092408435133414523352839687903472138975
\]
is the rational lower endpoint of the certified interval for \(\lambda_8\).

The cover contains 374 exactly adjacent intervals and no unresolved region: 319 with \(c_{512}>0\), 10 with \(c_{1024}<0\), two with \(c_{2048}>0\), all opposite to the sign of \(\tau\), and 43 near centers of periods 256, 512 or 1024, excluded by their second return.

In the latter case, \(g=f_\lambda^p\), \(d=g(0)\) and a uniform bound \(|g''|\le M\) between zero and \(d\) satisfy \(M|d|<2\). Then
\[
|g(d)-d|\le\tfrac12M|d|^2<|d|\quad(d\ne0).
\]
Both returns have the same sign; for \(d=0\), both are zero. Both situations exclude \(\tau_{2p}=-\tau_p\) with strict signs.

The evaluation uses rational moments, geometric tails of the potential and its second derivative, and \(|u_0'''|<6\). Parameter centering intersects the natural enclosure with
\[
c_j(\lambda_{\rm medio})+
\partial_\lambda c_j(I)\,(I-\lambda_{\rm medio});
\]
the derivative is propagated over the full box. The self-map property of \([-1,1]\) is proved before restricting the enclosures to it.

The complete certificate retains the cover and checks its exact tiling, together with every sign or Taylor inequality.

\subsubsection{Joint localization}
\label{hmtfg:escalado:15.4}

The local crossing \(\lambda_*\) realizes \(\tau\) through its real returns at every depth, so that \(a_*\le\lambda_*\). By \eqref{hmtfg:escalado:eq:50}, \(\lambda_8<a_*\); together with \eqref{hmtfg:escalado:eq:51},
\begin{equation}
\boxed{1.874038<a_*\le\lambda_*<1.874039.}
\label{hmtfg:escalado:eq:52}
\end{equation}
The localization preserves the canonical itinerary and places its first limiting parameter in the same interval as the stable crossing.

\subsection{First-exit isolation and coincidence of the limits}
\label{hmtfg:escalado:16}

This section completes the isolation identified in section~\ref{hmtfg:escalado:15.4}. It preserves the HMT family \eqref{hmtfg:escalado:eq:1}, the minimal selection \eqref{hmtfg:escalado:eq:50}, the chart \eqref{hmtfg:escalado:eq:5} and the stable sheet of sections~\ref{hmtfg:escalado:5}--\ref{hmtfg:escalado:7}. The proof uses the order \(a_*\le\lambda_*\): it suffices to exclude a separation toward the negative side of the expanding cone.

\subsubsection{Upper and lower bounds for transverse transport}
\label{hmtfg:escalado:16.1}

Besides \eqref{hmtfg:escalado:eq:9}, Lemma 2.1 of \href{https://arxiv.org/pdf/1410.3277}{Hertling–Spandl, section 2} provides \(\|D\Phi\|<0.9\) in the same ball of radius \(0.01\), where
\[
\Phi_u(u,\nu)=u-\frac{U(u,\nu)-u}{3.669}.
\]
Therefore,
\begin{equation}
\left|1-\frac{U_u-1}{3.669}\right|<0.9,\qquad
U_u<1+1.9(3.669)=7.9711.
\label{hmtfg:escalado:eq:53}
\end{equation}
The number \(3.669\) belongs to the published preconditioner of this analytic certificate; the generated family \eqref{hmtfg:escalado:eq:1} retains its prior coefficients.

Take \(\kappa=0.21\). For two points in the ball joined by a secant with
\(\Delta u<0\), \(\|\Delta\nu\|_1\le\kappa|\Delta u|\), integration of the blocks of \(DR\) along the segment gives
\begin{equation}
4.4034|\Delta u|
\le|\Delta u'|\le8.0887|\Delta u|,\qquad \Delta u'<0,
\label{hmtfg:escalado:eq:54}
\end{equation}
\begin{equation}
\|\Delta\nu'\|_1
\le(0.756+0.719\kappa)|\Delta u|
=0.90699|\Delta u|
<\kappa(4.4034)|\Delta u|.
\label{hmtfg:escalado:eq:55}
\end{equation}
Here \(4.4034=4.521-0.560\kappa\) and \(8.0887=7.9711+0.560\kappa\). Thus, the secant cone preserves orientation and expands while the segment remains in the ball.

\subsubsection{Reduction of an infinite separation to a finite region}
\label{hmtfg:escalado:16.2}

Suppose \(a_*<\lambda_*\), and set
\[
f_n=R^{4+n}f_{a_*},\qquad
s_n=R^{4+n}f_{\lambda_*},\qquad
(\Delta u_n,\Delta\nu_n)=f_n-s_n.
\]
The positive bound for \(u'\) and \(\|\nu'\|_1/u'<0.144386<\kappa\) on \(\Gamma(J)\) place the initial secant in the negative cone. The stable sheet is written relative to \(g\), with
\begin{equation}
s_n=g+e_n,\quad
\|e_n\|_{\mathcal A}<0.001666(0.83996)^n,\quad
|(e_n)_u|\le0.16\|(e_n)_\nu\|_1.
\label{hmtfg:escalado:eq:56}
\end{equation}
The entry checks imply
\begin{equation}
|\Delta u_0|
<0.004993128+0.00004
 +0.16(0.001396077+0.00004)
=0.00526290032< h,\qquad h=0.006.
\label{hmtfg:escalado:eq:57}
\end{equation}
While \(|\Delta u_n|<h\), \eqref{hmtfg:escalado:eq:55}–\eqref{hmtfg:escalado:eq:56} give
\begin{equation}
\|f_n-g_0\|_{\mathcal A}
<(1+\kappa)h+0.001666+0.00004
=0.008966<0.01.
\label{hmtfg:escalado:eq:58}
\end{equation}
The point \(s_n\) and the segment between the two also remain in that convex ball. The lower expansion \eqref{hmtfg:escalado:eq:54} forces a finite first exit \(N\), and the upper bound applied to the preceding step provides
\begin{equation}
0.006\le-\Delta u_N<0.0485322,\qquad
\|\Delta\nu_N\|_1\le0.21|\Delta u_N|.
\label{hmtfg:escalado:eq:59}
\end{equation}
The inequality \(8.0887h=0.0485322\) is used; the full jump, not only the first-exit surface, is included.

Consequently, \(f_N\) belongs to the following region:
\begin{equation}
\begin{split}
\mathcal C_-=\{\,g+(d_u,d_\nu)+e:\;&
-0.0485322\le d_u\le-0.006,\\
&\|d_\nu\|_1\le0.21|d_u|,\quad
|e_u|+\|e_\nu\|_1\le0.001666,\\
&|e_u|\le0.16\|e_\nu\|_1\,\}.
\end{split}
\label{hmtfg:escalado:eq:60}
\end{equation}
Since \(a_*\) realizes \(\tau\), all its normalized returns are unimodal self-maps of \([-1,1]\) with the same itinerary. This property follows from return induction, independently of the distance from \(f_N\) to \(g_0\). This is why the following certificate treats \(\mathcal C_-\) intersected with that class of self-maps.

\subsubsection{Certified exclusion of the first-exit region}
\label{hmtfg:escalado:16.3}

The rational first-exit exclusion certificate, reproduced in the computational supplement, covers the interval of \eqref{hmtfg:escalado:eq:59} exactly with 64 adjacent rational bands. All are excluded; no additional subdivision or unresolved region enters the result.

Of the 64 bands, 54 display an opposite sign and ten a contracting return of period 16 or 32. The largest Lipschitz bound for those returns is less than \(0.475146491\); the maximum exclusion horizon is iterate 64.

The evaluation retains the first forty perturbation coefficients as shared variables, together with an explicit bound on their infinite tail. For each band, the center is \(g_0\) shifted in coordinate \(u\). If \(a,b\ge0\) bound the dual sensitivities in \(u,\nu\), the support of the uncertainty in the fixed point and the stable residual is
\begin{equation}
0.00004\max(a,b)
+0.001666\max\left(b,\frac{0.16a+b}{1.16}\right).
\label{hmtfg:escalado:eq:61}
\end{equation}
The second expression is obtained by maximizing \(a|e_u|+b\|e_\nu\|_1\) under the two constraints of \eqref{hmtfg:escalado:eq:56}. This uses the orientation of the stable sheet proved previously, as well as its norm.

Let \(x_j\) be the center’s orbit and \(L_j\) its linear sensitivity to all shared coefficients. The Taylor model has the form
\[
c_j=x_j+L_j(\Delta v)+r_j,\qquad |r_j|\le E_j.
\]
If \(\rho_j\) bounds \(|L_j(\Delta v)|+E_j\), the outward recurrence used is
\begin{equation}
E_{j+1}\le |f_0'(x_j)|E_j
+\tfrac12\sup|f_0''|\rho_j^2
+\sup|\Delta f'|\rho_j.
\label{hmtfg:escalado:eq:62}
\end{equation}
The derivatives are bounded along the segment between the two states. The central orbit is verified within \([-1,1]\); the other orbits belong to that interval by the self-map property included in the certified domain. The calculations use outward rational rounding, including for the geometric tails.

Each band is excluded by one of two inequalities:

\par\noindent\textbf{1.}\enspace A critical value has a sign strictly opposite to \(\tau_j\).
\par\noindent\textbf{2.}\enspace For \(p=2^k\), the return \(H=f^p\) has Lipschitz constant \(L_p<1\) on the segment between \(0\) and \(H(0)\). Then
\[
|H(H(0))-H(0)|\le L_p|H(0)|.
\]
If \(H(0)\ne0\), the two values have the same sign; if \(H(0)=0\), both vanish. In both cases the strict itinerary \(\tau_{2p}=-\tau_p\) is excluded.

The constant \(L_p\) is obtained by propagating the initial segment \([-r,r]\), \(r\ge|H(0)|\), alongside the critical-orbit boxes and multiplying the outward derivative bounds. The maps of the full band are bounded, retaining \eqref{hmtfg:escalado:eq:61}.

\subsubsection{Coincidence theorem for the accumulation parameter}
\label{hmtfg:escalado:16.4}

\textbf{Theorem.} For the uniform HMT family \eqref{hmtfg:escalado:eq:1}, the limit of the original first-root selection coincides with the certified stable crossing:
\begin{equation}
\boxed{\lambda_n\nearrow a_*=\lambda_*,
\qquad 1.874038<\lambda_*<1.874039.}
\label{hmtfg:escalado:eq:63}
\end{equation}

\textbf{Proof.} \eqref{hmtfg:escalado:eq:50}–\eqref{hmtfg:escalado:eq:52} give \(\lambda_n\uparrow a_*\le\lambda_*\), with both parameters in \(J\). If the inequality were strict, \eqref{hmtfg:escalado:eq:54}–\eqref{hmtfg:escalado:eq:59} would produce a return \(f_N\in\mathcal C_-\) realizing \(\tau\). The complete cover of section~\ref{hmtfg:escalado:16.3} excludes precisely that possibility. Hence \(a_*=\lambda_*\).

The proof controls every depth through a finite first exit and a uniform cover of its region. The identity \eqref{hmtfg:escalado:eq:63} links the global selector’s limit to the transverse crossing.


\subsection{Identification of the original centers and global scaling}
\label{hmtfg:escalado:17}

The final step preserves the selection of \eqref{hmtfg:escalado:eq:50}. The local continuation and the original parameters are compared through their exact period and their hybrid class; coincidence is obtained through uniqueness in a fixed parameter neighborhood.

\subsubsection{Degree-two analytic realization and transversality}
\label{hmtfg:escalado:17.1}

The fixed point \(g\) of section~\ref{hmtfg:escalado:3} is even, analytic, with a quadratic critical point and stationary doubling combinatorics. Theorem 2.1 of \href{https://annals.math.princeton.edu/wp-content/uploads/annals-v164-n3-p01.pdf}{de Faria–de Melo–Pinto}, applied to that single combinatorics, gives a one-element limit set, with a proper holomorphic extension of degree two between nested topological disks. Its attraction on the interval identifies that element with \(g\), since \(Rg=g\).

Lemma 3.2 of the same source extends a fixed number of returns to a full analytic neighborhood of \(g\), with degree-two images and positive modulus. The norm \eqref{hmtfg:escalado:eq:5} controls the uniform norm on a sufficiently small complex neighborhood of \([-1,1]\), because \(|(x^2-1)/2.5|<1\) there. The convergence \eqref{hmtfg:escalado:eq:3} enters that neighborhood; analytic dependence of the returns provides, for some fixed integer \(d\), a family
\begin{equation}
\mathcal F_\lambda=R^d f_\lambda
\label{hmtfg:escalado:eq:64}
\end{equation}
of degree-two maps, analytic in a complex neighborhood of \(\lambda_*\). The disks and the number of returns are fixed before varying \(\lambda\) in that neighborhood.

The parameter tangent belongs to the expanding cone, whereas the stable sheet satisfies the slope bound of section~\ref{hmtfg:escalado:5}. The returns preserve that separation. The degree-two realization locally identifies the stable sheet with the hybrid class of the infinitely renormalizable map; hence \eqref{hmtfg:escalado:eq:64} is transverse to that class.

The change of analytic space is justified in detail in \ref{hmtfg:clasificacion:8.2}, devoted to the transport of transversality. A direction tangent to the hybrid class would produce decreasing renormalized tangents, by uniform convergence on a hybrid disk and Cauchy’s estimate. The certified direction, by contrast, satisfies \(|\dot f(1)|=|\dot u|/10\) and grows under the returns by the expanding cone. This common evaluation of both representatives proves transversality without assuming equivalence of their norms.

\subsubsection{Uniqueness in a fixed neighborhood}
\label{hmtfg:escalado:17.2}

For \(n>d\), the classification of section~\ref{hmtfg:escalado:15} gives
\begin{equation}
K(\mathcal F_{\lambda_n})=W_{n-d}0.
\label{hmtfg:escalado:eq:65}
\end{equation}
Straightening a degree-two map preserves the critical orbit; \eqref{hmtfg:escalado:eq:65} identifies its hybrid class with that of the canonical quadratic center of period \(2^{n-d}\). The a priori bounds for the real doubling sequence are those used in Lyubich’s universality theorem.

\href{https://www.maths.tcd.ie/EMIS/journals/Annals/149_2/lyubich.pdf\#page=69}{Lyubich’s Theorem 7.4, pp. 387–388} gives a single intersection with each of those classes, at sufficiently high levels, on the analytic transversal and within a fixed neighborhood of the accumulation parameter.

By \eqref{hmtfg:escalado:eq:63}, \(\lambda_n\) eventually enters that neighborhood. The local continuation of section~\ref{hmtfg:escalado:8}, indexed by the same period, belongs to the same hybrid class. Uniqueness implies
\begin{equation}
\boxed{\lambda_n=\widehat\mu_n\quad(n\ge n_0),}
\label{hmtfg:escalado:eq:66}
\end{equation}
where \(\widehat\mu_n\) denotes the local center of physical period \(2^n\). This is a reindexing of \(\mu_j\) by a fixed shift determined by the returns and the target of section~\ref{hmtfg:escalado:8}. The original selection of \(\lambda_n\) remains intact.

The integer \(n_0\) is existential in this proof. Continuation at every level and the word \(W_n0\) are already proved in \eqref{hmtfg:escalado:eq:50}; \eqref{hmtfg:escalado:eq:66} provides the metric identity needed for the limit.

\subsubsection{Scaling theorem for the first HMT roots}
\label{hmtfg:escalado:17.3}

\textbf{Theorem.} For the uniform family \eqref{hmtfg:escalado:eq:1}, let \(\lambda_n\) be the successive first new roots of exact critical period \(2^n\). There exist \(C>0\), \(0<\theta<1\) and \(M<\infty\) such that, for every sufficiently large \(n\),
\begin{equation}
\lambda_*-\lambda_n=C\delta_{\mathrm F}^{-n}(1+e_n),
\qquad |e_n|\le M\theta^n,
\label{hmtfg:escalado:eq:67}
\end{equation}
and
\begin{equation}
\boxed{\lim_{n\to\infty}
\frac{\lambda_{n-1}-\lambda_{n-2}}
{\lambda_n-\lambda_{n-1}}=\delta_{\mathrm F}.}
\label{hmtfg:escalado:eq:68}
\end{equation}

\textbf{Proof.} \eqref{hmtfg:escalado:eq:66} transports the local scaling \eqref{hmtfg:escalado:eq:24} to the original roots; the finite index shift is absorbed into \(C,M\). Its sign is positive because \(\lambda_n<\lambda_*\).

The power-law remainder also admits a direct justification through the holonomy of hybrid classes. Lyubich’s Lemma 7.3 gives \(C^{1+\beta}\) regularity, with nonzero derivative, on the transversal at the accumulation point. The one-dimensional unstable dynamics is analytic and has multiplier \(\delta_{\mathrm F}>1\); its linearizing coordinate places successive centers at distances proportional to \(\delta_{\mathrm F}^{-n}\). Composing with the inverse holonomy gives a nonzero linear term and a remainder \(O(\delta_{\mathrm F}^{-(1+\beta)n})\). Thus one may take \(\theta=\delta_{\mathrm F}^{-\beta}\).

Writing
\[
r_n=\frac{\delta_{\mathrm F} e_{n-1}-e_n}{\delta_{\mathrm F}-1},
\qquad K=\frac{M(\delta_{\mathrm F}+\theta)}{\delta_{\mathrm F}-1},
\]
one obtains
\[
\lambda_n-\lambda_{n-1}
=C(\delta_{\mathrm F}-1)\delta_{\mathrm F}^{-n}(1+r_n),\qquad
|r_n|\le K\theta^{n-1}.
\]
Therefore,
\begin{equation}
\left|
\frac{\lambda_{n-1}-\lambda_{n-2}}
{\lambda_n-\lambda_{n-1}}-\delta_{\mathrm F}
\right|
\le
\frac{\delta_{\mathrm F} K(1+\theta)\theta^{n-2}}
{1-K\theta^{n-1}},
\quad K\theta^{n-1}<1.
\label{hmtfg:escalado:eq:69}
\end{equation}
The right-hand side tends to zero, proving \eqref{hmtfg:escalado:eq:68}.

\subsubsection{Chain of results and scope of the closure}
\label{hmtfg:escalado:17.4}

The effective composition is: APP–TRIT–TPK dodecaphase reader; profile \eqref{hmtfg:escalado:eq:1}; ordered transport of returns \eqref{hmtfg:escalado:eq:38}–\eqref{hmtfg:escalado:eq:40}; classification and existence of the first roots \eqref{hmtfg:escalado:eq:50}; localization \eqref{hmtfg:escalado:eq:52}; first-exit isolation \eqref{hmtfg:escalado:eq:63}; eventual identity of centers \eqref{hmtfg:escalado:eq:66}; scaling \eqref{hmtfg:escalado:eq:67}–\eqref{hmtfg:escalado:eq:69}. The Jacobi representation \eqref{hmtfg:escalado:eq:43} preserves this reader’s full profile and its marks remain in the antecedent.

The spatial expansion of the HMT reader and nonadic refinement provide the genealogy and evaluation at arbitrary depth. Scaling between parameters is proved through the return dynamics, transversality and isolation that we have just composed. The cited analytic proofs act on the family already constructed; their comparison values do not select its phases or coefficients.

The closure concerns the original global selector of the uniform sector \(\eta=0\), with the listed theorem dependencies. The asymptotic law retains its existential quantifiers. The width of a finite-root box measures evaluation precision; the estimate \eqref{hmtfg:escalado:eq:69} expresses convergence toward the limit.


The programs and certificates in the reproducibility supplement make it possible to reproduce the finite parts of the proof. Continuation at every depth rests on the inductions, the uniform first-exit argument and the identification theorems, with their explicit domains.
